\documentclass[a4paper,10pt]{extarticle}

% Packages
\usepackage[utf8]{inputenc}
\usepackage{geometry}
\geometry{a4paper, margin=0.4in}
\usepackage{titlesec}
\usepackage{enumitem}
\usepackage{hyperref} % For clickable links
\usepackage{fontawesome5} % FontAwesome for icons
\usepackage{polski}

% Formatting
\setlist{noitemsep}
\titleformat{\section}{\large\bfseries}{\thesection}{1em}{}[\titlerule]
\titlespacing*{\section}{0pt}{\baselineskip}{\baselineskip}

% Begin document
\begin{document}

% Disable page numbers
\pagestyle{empty}

% Header
\begin{center}
    \fontsize{26 pt}{26 pt}\selectfont
    \textbf{Daniel Ambroziewicz}

    \vspace{0.3 cm}

    \normalsize
    \mbox{{\footnotesize\faMapMarker*}\hspace*{0.13cm}Kraków, Polska}% 
    \kern 0.25 cm%
    \mbox{\href{mailto:ambroziewiczdaniel@gmail.com}{{\footnotesize\faEnvelope[regular]}\hspace*{0.13cm}ambroziewiczdaniel@gmail.com}}%
    \kern 0.25 cm%
    \mbox{\href{tel:+48-512-337-500}{{\footnotesize\faPhone*}\hspace*{0.13cm}512 337 500}}%
    \kern 0.25 cm%
    \mbox{\href{https://www.linkedin.com/in/daniel-ambroziewicz-9343a414a/}{{\footnotesize\faLinkedinIn}\hspace*{0.13cm}Daniel Ambroziewicz}}%
    \kern 0.25 cm%
\end{center}

\section*{DOŚWIADCZENIE ZAWODOWE}

\noindent
\textbf{Process Innovation Intern} w \textbf{Aptiv (Versigent)} \hfill Czerwiec 2025 – Obecnie
\begin{itemize}
    \item Implementacja i rozwój systemów wizji komputerowej przy użyciu języka Python, biblioteki PyTorch, algorytmów YOLO oraz klasycznych metod wizyjnych w celu wsparcia innowacji sprzętu produkcyjnego
    \item Współtwórca innowacyjnego rozwiązania technologicznego, docenionego i chronionego wewnętrznie jako tajemnica przedsiębiorstwa (Trade Secret)
    \item Projektowanie i integracja systemów wbudowanych (m.in. \textbf{Arduino}) z układami mechanicznymi, sterowanych za pomocą języka \textbf{Python} w ramach optymalizacji i innowacji procesów produkcyjnych
    \item Inżynieria wsteczna i adaptacja starszej aplikacji w języku C\# zapewniająca kompatybilność z oprogramowaniem BricsCAD oraz integrację z istniejącymi przepływami pracy CAD
    \item Tworzenie licznych aplikacji w języku Python, wykorzystywanych obecnie w produkcji wiązek elektrycznych
\end{itemize}

\vspace{0.15cm}
\noindent
\textbf{Kontroler jakości} w \textbf{Superior Industries} \hfill Maj 2022 – Sierpień 2022
\begin{itemize}
    \item Kontrola wyrobów pod kątem uszkodzeń mechanicznych oraz praca zespołowa na hali produkcyjnej
\end{itemize}

\section*{WYKSZTAŁCENIE}
\noindent
\textbf{Akademia Górniczo-Hutnicza im. Stanisława Staszica w Krakowie} \\
Wydział Inżynierii Mechanicznej i Robotyki 
\begin{itemize}
    \item \textbf{Studia magisterskie: Inżynieria mechatroniczna} \hfill Luty 2026 – Obecnie \\
    Specjalizacja: \textit{Projektowanie Mechatroniczne - Informatyka} (1. semestr)
    \item \textbf{Studia inżynierskie: Inżynieria mechatroniczna} \hfill Październik 2022 – Luty 2026 \\
    Tytuł: Inżynier
\end{itemize}

\section*{PROJEKTY}

\noindent
\textbf{Zautomatyzowany system wizyjny kontroli jakości (Praca inżynierska)} \hfill 2025 – 2026
\begin{itemize}
    \item Samodzielne zaprojektowanie i budowa stanowiska wyposażonego w przenośnik taśmowy (sterowanie \textbf{Arduino}) oraz zintegrowany system wizyjny
    \item Opracowanie w języku \textbf{Python} hybrydowego algorytmu detekcji wad, łączącego klasyczne metody przetwarzania obrazu (\textbf{OpenCV}) z sieciami neuronowymi (\textbf{YOLOv8, Anomalib, PyTorch})
    \item System osiągnął dokładność klasyfikacji defektów na poziomie ponad 97\%
    \item \textbf{Wyróżnienie:} Projekt doceniony przez władze uczelni (oficjalne wideo promocyjne AGH) oraz opisany w prasie (m.in. \href{https://gazetakrakowska.pl/student-krakowskiej-agh-zaprojektowal-system-do-zautomatyzowanej-kontroli-jakosci-ciastek/gh/c1p2-28735765}{\textit{Gazeta Krakowska}} i \href{https://dziennikpolski24.pl/student-krakowskiej-agh-zaprojektowal-system-do-zautomatyzowanej-kontroli-jakosci-ciastek/ar/c1p2-28735765}{\textit{Dziennik Polski}})
\end{itemize}

\vspace{0.15cm}
\noindent
\textbf{{Projekt i budowa prototypu robota mobilnego}} \hfill 2023
\begin{itemize}
    \item Konstrukcja robota z wysuwanym mostem do pokonywania przeszkód (projekt 3-osobowy)
    \item Osobiste wykonanie algorytmu sterowania w języku \textbf{C++} (ocena 100/100 za część programistyczną)
\end{itemize}

\section*{UMIEJĘTNOŚCI}

\noindent
\textbf{{Programowanie i IT}} \hfill
\begin{itemize}
    \item Języki: \textbf{Python, C++, C, C\#}, MatLab, HTML, CSS, LaTeX
    \item Sztuczna inteligencja i widzenie maszynowe: \textbf{PyTorch, YOLO, Computer Vision}
\end{itemize}

\noindent
\textbf{{Oprogramowanie inżynierskie i symulacje}} \hfill
\begin{itemize}
    \item Pakiet CAD/CAM: \textbf{AutoCAD, Autodesk Inventor, SolidWorks, BricsCAD}
    \item Analiza i MES: \textbf{Siemens NX, MSC Nastran, Hexagon/MSC.Software}
    \item Analiza danych: \textbf{MatLab, Simulink} (w tym algorytmy sieci neuronowych i logiki rozmytej)
\end{itemize}

\noindent
\textbf{{Robotyka i CNC}} \hfill
\begin{itemize}
    \item Programowanie robotów: \textbf{Omron} (język V++, MobilePlanner), \textbf{Mitsubishi} (MELFA BASIC IV), \textbf{Boston Dynamics} (Python)
    \item Oprogramowanie off-line: \textbf{RobotStudio, ACE}
    \item Obróbka maszynowa (CNC): \textbf{G-Code, HEIDENHAIN}
\end{itemize}

\noindent
\textbf{{Elektronika i Systemy Wbudowane}} \hfill 
\begin{itemize}
    \item Systemy wbudowane: \textbf{Arduino}, programowanie układów \textbf{FPGA (NIOS)} przy użyciu \textbf{Intel Quartus Prime} i języka C
    \item Projektowanie obwodów elektronicznych (fizycznie i wirtualnie - \textbf{MultiSim})
\end{itemize}


\section*{JĘZYKI}
\noindent
Angielski (B2 - certyfikat SJO AGH), Polski (język ojczysty)

\vfill
\renewcommand{\thefootnote}{}
\footnotetext{\fontsize{8 pt}{8 pt}\selectfont
Wyrażam zgodę na przetwarzanie moich danych osobowych dla potrzeb niezbędnych do realizacji procesu rekrutacji zgodnie z Rozporządzeniem Parlamentu Europejskiego i Rady (UE) 2016/679 z dnia 27 kwietnia 2016 r. w sprawie ochrony osób fizycznych w związku z przetwarzaniem danych osobowych i w sprawie swobodnego przepływu takich danych oraz uchylenia dyrektywy 95/46/WE (RODO).}

\end{document}